% =====================================================================
%  main.tex — MASTER FILE (compile this one)
%
%  A simplified, verified guide to HyperPot: a fault-tolerant distributed
%  LLM inference system that pools consumer laptops over Wi-Fi. Explains
%  the idea, WHY it is defensible (with evidence vs. existing projects),
%  the exact technology stack, the hardware, the MVP plan, benchmarks,
%  and the Meta Global AI Developer Hackathon (Indie track) submission
%  strategy. Every important claim is tagged by evidence strength:
%  [VERIFIED] fact from a source, [HYPOTHESIS] reasoned, [UNKNOWN] unproven.
%
%  Build:  pdflatex main   (run twice for the table of contents)
%  Each section is its own file under sections/ so it is easy to edit.
% =====================================================================
\documentclass[11pt]{article}
% =====================================================================
%  preamble.tex — shared setup for the Simplified HyperPot guide
%  Edit formatting/macros here; edit prose in sections/*.tex.
% =====================================================================
\usepackage[T1]{fontenc}
\usepackage[utf8]{inputenc}
\usepackage[a4paper,margin=1in]{geometry}
\usepackage{amsmath,amssymb}
\usepackage{booktabs}
\usepackage{tabularx}
\usepackage{array}
\usepackage{enumitem}
\usepackage{xcolor}
\usepackage{tikz}
\usepackage{hyperref}
\usepackage{listings}
\usepackage{microtype}
\usepackage{fancyhdr}
\usepackage{textcomp}
\usepackage{newunicodechar}
\usetikzlibrary{arrows.meta,positioning,shapes.geometric,calc,fit,backgrounds}

\extrafloats{120}

% Make Unicode symbols used in the prose safe under pdfLaTeX.
\newunicodechar{₹}{Rs.\,}
\newunicodechar{→}{\ensuremath{\rightarrow}}
\newunicodechar{≤}{\ensuremath{\leq}}
\newunicodechar{≥}{\ensuremath{\geq}}
\newunicodechar{≠}{\ensuremath{\neq}}
\newunicodechar{×}{\ensuremath{\times}}
\newunicodechar{§}{\S}
\newunicodechar{–}{--}
\newunicodechar{—}{---}
\newunicodechar{✓}{\ensuremath{\checkmark}}

\hypersetup{colorlinks=true,linkcolor=blue!55!black,urlcolor=blue!55!black,citecolor=blue!55!black}

% ---- Column types -----------------------------------------------------
\newcolumntype{Y}{>{\raggedright\arraybackslash}X}
\newcolumntype{L}[1]{>{\raggedright\arraybackslash}p{#1}}
\newcolumntype{R}[1]{>{\raggedleft\arraybackslash}p{#1}}

% ---- Code / route helpers --------------------------------------------
\newcommand{\code}[1]{\texttt{\small #1}}
\newcommand{\route}[1]{\par\noindent{\small\colorbox{gray!12}{\parbox{0.97\linewidth}{\ttfamily #1}}}\par\smallskip}

% ---- Evidence tags (the HyperPot equivalent of data-source tags) -----
\definecolor{vercol}{RGB}{20,110,60}
\definecolor{hypcol}{RGB}{150,90,0}
\definecolor{unkcol}{RGB}{150,40,40}
\newcommand{\VER}{\textcolor{vercol}{\textbf{[VERIFIED]}}}
\newcommand{\HYP}{\textcolor{hypcol}{\textbf{[HYPOTHESIS]}}}
\newcommand{\UNK}{\textcolor{unkcol}{\textbf{[UNKNOWN]}}}

% ---- Keep / Drop tags -------------------------------------------------
\definecolor{keepcol}{RGB}{22,120,60}
\definecolor{dropcol}{RGB}{170,30,30}
\newcommand{\Keep}{\textcolor{keepcol}{\textbf{KEEP}}}
\newcommand{\Drop}{\textcolor{dropcol}{\textbf{DROP}}}

% ---- Reusable "field" lead-in -----------------------------------------
\newcommand{\field}[1]{\par\smallskip\noindent\textbf{#1.}\ }

% ---- Listings ---------------------------------------------------------
\lstdefinestyle{plain}{basicstyle=\ttfamily\footnotesize,backgroundcolor=\color{gray!8},
  frame=single,rulecolor=\color{gray!35},breaklines=true,columns=fullflexible,showstringspaces=false}

% ---- Headers ----------------------------------------------------------
\pagestyle{fancy}
\fancyhf{}
\fancyhead[L]{\small HyperPot — Simplified Guide}
\fancyhead[R]{\small\thepage}
\renewcommand{\headrulewidth}{0.4pt}


\title{\vspace{-1.5em}\textbf{HyperPot}\\[0.2em]
\large Fault-Tolerant Distributed LLM Inference on a Pool of Consumer Laptops\\[0.3em]
\normalsize A Simplified, Evidence-Tagged Idea \& Build Guide}
\author{\normalsize HyperPot --- Meta Global AI Developer Hackathon (Indie Track)}
\date{\normalsize Version 1.0}

\begin{document}
\maketitle
\thispagestyle{fancy}

\begin{abstract}
\noindent This document explains HyperPot in plain language: what it is, the one
problem it solves, and exactly how it will be built. HyperPot pools ordinary
student laptops --- across different GPUs and operating systems, connected over
Wi-Fi --- into a single cluster that runs a Meta Llama model. Its defining
capability is \textbf{fault tolerance}: when one laptop drops off the network
mid-generation (lid closed, Wi-Fi lost, battery dead), the system detects it,
re-routes that laptop's share of the work to a surviving machine, and the answer
keeps streaming. Every important claim here is tagged \VER\ (verified from a
source), \HYP\ (a reasoned hypothesis), or \UNK\ (unknown / unproven), so the
reader knows exactly what to trust. The guide is honest about novelty: the
\emph{algorithm} of load-aware failure recovery already exists in datacenter
research; HyperPot's contribution is bringing it down to a pile of cheap,
unreliable laptops on Wi-Fi using only free tools --- and making it demonstrable
live. The final chapters give the technology stack, the hardware map, a
phase-by-phase MVP plan to the submission deadline, measurable benchmarks, and
the hackathon strategy.
\end{abstract}

\tableofcontents
\clearpage

% ---------------------------------------------------------------------
\section{How to Read This Guide}
% ---------------------------------------------------------------------

HyperPot turns a collection of ordinary laptops into one shared computer that
runs a large language model (LLM). This guide documents the \emph{idea}, the
\emph{evidence} behind it, and the exact \emph{plan} to build it for the Meta
Global AI Developer Hackathon (Indie track). It is written so a newcomer can
follow it, while remaining strictly honest about what is proven and what is not.

\subsection{The evidence tags}
Next to important claims you will see one of three tags. They answer the
question ``how much should I trust this?'':

\begin{itemize}[leftmargin=1.4em,nosep]
  \item \VER\ --- a \textbf{verified fact} taken from an official repository,
  research paper, or the hackathon's own rules. These can be stated to a judge.
  \item \HYP\ --- a \textbf{reasoned hypothesis}: a sensible inference that has
  not yet been proven by a test or a source.
  \item \UNK\ --- \textbf{unknown}: not enough evidence yet; must be checked
  before relying on it.
\end{itemize}

\subsection{Keep vs drop}
Design decisions are marked \Keep\ (part of the plan) or \Drop\ (deliberately
excluded, with a reason). Being explicit about what we are \emph{not} building
is as important as what we are.

\subsection{The one-paragraph mental model}
A single LLM is split into layers. Each laptop (a ``worker'') loads a slice of
those layers. To answer a prompt, data flows worker~1 → worker~2 → worker~3 like
an assembly line, and words come out the end. A central ``controller'' hands out
the work, listens for heartbeats, and --- the key part --- when a worker goes
silent, it reassigns that worker's slice to a healthy machine and re-runs the
request so the answer continues. The whole point is that the laptops are
\emph{unreliable by nature} (they are people's real machines on Wi-Fi), and
HyperPot makes the cluster behave reliably anyway.

\subsection{The honest novelty statement}
We do not claim to have invented load-aware failure recovery for LLMs. That
already exists in datacenter research (\S\ref{sec:landscape}). HyperPot's
contribution, stated safely, is: \emph{``Fault tolerance like this exists in
datacenters. It has never run on a pile of cheap student laptops over Wi-Fi with
free tools. We made that work --- and here, watch, we kill one.''} This framing
is defensible because every part of it is backed by a source in
\S\ref{sec:landscape}.

% ---------------------------------------------------------------------
\section{The Idea}
\label{sec:idea}
% ---------------------------------------------------------------------

\subsection{What HyperPot is}
HyperPot is a distributed LLM system that uses \textbf{consumer laptops as shared
computing resources}. Instead of renting a cloud GPU or needing one expensive
machine, several everyday laptops pool their compute and memory to run a model
together.

\subsection{The original architecture (the starting point)}
The project begins from an existing three-part design:

\begin{itemize}[leftmargin=1.4em,nosep]
  \item \textbf{Go} --- the master controller: HTTP API, task queue, worker
  heartbeats, and result collection. \Keep
  \item \textbf{Rust} --- the worker agents that connect to the controller and
  execute inference tasks. \Keep
  \item \textbf{llama.cpp} --- the LLM inference engine (originally prebuilt
  Windows binaries). \Keep, but used as a \emph{local} engine only (see
  \S\ref{sec:techstack}).
\end{itemize}

\subsection{The defining problem it solves}
The laptops belong to real people and sit on Wi-Fi, so they come and go: a lid
closes, Wi-Fi drops, a battery dies. HyperPot's defining capability is to
\textbf{survive that} --- to keep generating the answer even when a laptop
vanishes mid-sentence. This is the one thing the project is built around.

\subsection{What this guide deliberately is NOT about}
\begin{itemize}[leftmargin=1.4em,nosep]
  \item \Drop\ \textbf{Distributed fine-tuning / training.} It overlaps existing
  projects (Petals, QVAC Fabric), is far harder to make fault-tolerant, and adds
  nothing to the live demo. HyperPot is \textbf{inference only}. If personalization
  is wanted, a small LoRA adapter can be trained offline on one machine
  beforehand --- never distributed training during the event. \HYP
  \item \Drop\ \textbf{Beating anyone on raw speed.} Faster heterogeneous
  inference is already owned by prima.cpp (\S\ref{sec:landscape}). We do not
  compete there.
\end{itemize}

% ---------------------------------------------------------------------
\section{The Competitive Landscape (Verified)}
\label{sec:landscape}
% ---------------------------------------------------------------------

This chapter is the evidence base. Everything here was checked against official
repositories, papers, or issue trackers.

\subsection{The existing projects}
\begin{table}[h]
\centering\footnotesize
\caption{What existing distributed-inference projects actually do}
\label{tab:landscape}
\begin{tabularx}{\linewidth}{L{0.16\linewidth}Y L{0.26\linewidth}}
\toprule
Project & What it does & Fault tolerance on node loss \\
\midrule
llama.cpp RPC & Offloads layers to remote \code{rpc-server} nodes over TCP &
\textbf{None} --- a peer crash/drop ``terminates the process'' \VER \\
Exo & Auto-discovery cluster, best on Apple Silicon (MLX, Thunderbolt) &
``Critical gaps in fault tolerance'' per 2025 review + open issues \VER \\
Petals & BitTorrent-style swarm for 50B+ models over the Internet; also fine-tuning &
Designed for churn, but Internet-scale / large models (not LAN laptops) \VER \\
Distributed Llama & Home-device tensor parallel over Ethernet &
Static topology; no churn-recovery claim \VER \\
prima.cpp & 30--70B on heterogeneous home clusters; Halda scheduler; Wi-Fi tolerant &
Static, up-front scheduling; no mid-inference recovery emphasis \VER \\
QVAC Fabric & Edge-first inference + LoRA fine-tuning on phones/laptops/GPUs &
Marketing claims decentralization; internals \UNK \\
\bottomrule
\end{tabularx}
\end{table}

\subsection{The key verified facts}
\begin{itemize}[leftmargin=1.4em,nosep]
  \item \VER\ \textbf{llama.cpp RPC has no fault tolerance.} Its own docs state a
  peer crash, a dropped network, or a malformed response terminates the process
  (via \code{RPC\_STATUS\_ASSERT}/\code{GGML\_ABORT}). Confirmed by GitHub issues
  \#11508 and \#9214.
  \item \VER\ \textbf{Exo's fault tolerance is weak in practice.} A Sept 2025
  deep-dive review cites ``critical gaps in security, fault tolerance, and
  operational tooling''; open issues show nodes dropping and unstable node counts.
  \item \VER\ \textbf{Petals already solved churn-tolerant inference} --- but
  over the \emph{Internet} for \emph{50B+} models, not a LAN laptop pool.
  \item \VER\ \textbf{Capability-aware failure recovery already exists in
  research}: LUMEN (arXiv 2606.17787) does ``interrupted-request distribution at
  failure time'' considering cluster load; ``Intelligent Orchestration \ldots at
  the Edge'' (arXiv 2504.03668) does capacity-aware distribution + dynamic
  partition migration; FailSafe (arXiv 2511.14116) survives multiple GPU
  failures. These target \textbf{GPU/cloud/edge-server} clusters, not consumer
  laptops on Wi-Fi with llama.cpp.
\end{itemize}

\subsection{What this means for novelty}
\begin{itemize}[leftmargin=1.4em,nosep]
  \item As a \emph{research algorithm}, load-aware recovery is \textbf{not
  novel}. \VER\ Do not claim invention to a judge.
  \item As a \emph{working system} on consumer laptops, over Wi-Fi, using
  llama.cpp (which cannot do it) --- that combination is \textbf{unclaimed} as
  far as searched. This is \textbf{engineering / integration novelty} plus a
  strong live demo. \HYP\ (absence of a shipped product; not a proof of absence)
\end{itemize}

% ---------------------------------------------------------------------
\section{What Fault Tolerance Solves (In Plain Words)}
\label{sec:solves}
% ---------------------------------------------------------------------

\subsection{The group-essay analogy}
Imagine four friends writing one long essay together, split into parts. Halfway
through, one friend's laptop dies and he leaves.
\begin{itemize}[leftmargin=1.4em,nosep]
  \item \textbf{Without fault tolerance:} the whole essay is ruined; everyone
  starts over from the beginning. All prior work is wasted.
  \item \textbf{With fault tolerance:} the group notices he is gone, hands his
  part to a friend who has a spare copy of the material, and the essay continues.
  A small hiccup, then back to normal.
\end{itemize}
Fault tolerance = \textbf{the system keeps working even when a part of it breaks}.

\subsection{The LLM version}
HyperPot splits one model across laptops in an assembly line (worker~1 →
worker~2 → worker~3). Because these are real laptops on Wi-Fi, any one can
disappear at any moment. With today's tools (llama.cpp RPC) this is fatal: one
laptop disconnecting crashes the whole run and the half-written answer is lost
(\VER, \S\ref{sec:landscape}). HyperPot fixes exactly this: a laptop drops →
the controller notices via heartbeat → its share is re-routed to a surviving
laptop → the answer keeps streaming.

\subsection{A real-time example}
\field{Scene} A college lab: student laptops on the lab Wi-Fi together run a
Llama model to answer questions for a class. A student asks, ``Explain
photosynthesis step by step\ldots'' and the answer is half-printed on screen.
Suddenly one student \textbf{grabs their laptop and walks out} --- that laptop
held one quarter of the model.
\begin{itemize}[leftmargin=1.4em,nosep]
  \item \textbf{Today's systems:} the answer freezes and errors out mid-sentence.
  \item \textbf{HyperPot:} a \textasciitilde1--2 second pause, then the answer
  continues: ``\ldots the plant uses sunlight to turn water and CO\textsubscript{2}
  into sugar.'' The user never knew a laptop just left the building.
\end{itemize}

\subsection{The improved recovery idea (capability-aware)}
Rather than dumping the dead laptop's whole slice onto one machine, the
controller can \textbf{check who has spare compute} and split the slice by
capability (a strong RTX~4050 takes more layers than a weaker RTX~3050). This is
smarter, but note the tradeoff below.

\field{The engineering catch \HYP} A laptop can only take over a slice if it
already holds those model layers in memory. So capability-aware redistribution
requires survivors to \emph{pre-load overlapping slices}, which costs memory on
every laptop. Two options:
\begin{itemize}[leftmargin=1.4em,nosep]
  \item \textbf{Simple (recommended for MVP):} one identical hot-standby holds a
  copy of a worker's slice $\rightarrow$ fast, reliable recovery.
  \item \textbf{Smarter (stretch):} several laptops pre-hold overlapping slices
  $\rightarrow$ can rebalance by power, but uses more memory and is harder to
  finish in time.
\end{itemize}

\subsection{What it solves, in one line}
\begin{quote}
It makes a cluster of unreliable, come-and-go laptops behave like one reliable
computer --- so people's real laptops dropping off the Wi-Fi does not destroy
the work.
\end{quote}

% ---------------------------------------------------------------------
\section{The Technology Stack}
\label{sec:techstack}
% ---------------------------------------------------------------------

Everything here \emph{extends} the original Go~+~Rust~+~llama.cpp design; nothing
is rewritten. The one mindset shift: \textbf{the controller owns reliability,
not llama.cpp}.

\subsection{Component by component}

\field{Controller --- Go \Keep} The coordinator ``brain''. Handles the HTTP API
(user sends prompts), the task queue, worker heartbeats (detect a dead laptop),
result collection, and --- the novel part --- the \textbf{failover logic}:
``node died → who is alive and has capacity → reassign its layers → re-run the
request''. Go is well suited to this and the heartbeat/queue pieces already exist.

\field{Worker agent --- Rust \Keep} Runs on each laptop. Registers with the
controller, sends heartbeats (``alive; here is my CPU/RAM/battery''), loads its
assigned slice of the model, runs its layers via llama.cpp, and passes
activations to the next worker. Rust compiles to Windows, Linux, and macOS.

\field{Inference engine --- llama.cpp \Keep (usage changed)}
\begin{itemize}[leftmargin=1.4em,nosep]
  \item \Drop\ \textbf{Do NOT use llama.cpp's built-in RPC mode.} It crashes on
  node loss (\VER, \S\ref{sec:landscape}); building on it inherits the exact
  failure we are beating.
  \item \Keep\ \textbf{Use llama.cpp as a local engine on each worker}, with the
  Go controller in charge of coordination and recovery. We route \emph{around}
  llama.cpp's weak part. This single decision is what makes the project possible.
\end{itemize}

\field{Model --- Meta Llama, small, GGUF \Keep} A small open-weights Llama (e.g.
Llama~3.2 1B or 3B) in GGUF, quantized (e.g. Q4). Small = fast load, fits weak
laptops, snappy demo. Satisfies the hackathon's ``use a Meta model'' requirement.
\Drop\ 70B models --- slow demo, and prima.cpp wins that contest anyway.

\field{Networking --- TCP/gRPC over Wi-Fi \Keep} Workers $\leftrightarrow$
controller over plain TCP/JSON to start (upgrade to gRPC only if needed). Wi-Fi
is both the demo environment and the story (``real laptops, real Wi-Fi'').

\field{Model split --- pipeline (by layers) \Keep} Split by layers (worker~1 =
layers 0--8, worker~2 = 9--16, \ldots). For recovery, a survivor picks up the
dead worker's layers.

\field{Recovery method --- request-level re-route \Keep (the core)}
When a worker dies mid-request, the controller reassigns its layers to a survivor
and \textbf{re-runs the current request} from the start of generation. The user
sees a short pause, then continuation. \Drop\ mid-token activation handoff --- too
hard for the time window.

\field{Code written with --- Kiro \Keep} The Go controller, the Rust worker, the
failover handling, and the glue are written with Kiro.

\subsection{Summary table}
\begin{table}[h]
\centering\footnotesize
\caption{HyperPot technology stack}
\label{tab:stack}
\begin{tabularx}{\linewidth}{L{0.34\linewidth}Y}
\toprule
Piece & Choice \\
\midrule
Controller / failover brain & Go (keep, extend) \\
Worker agent per laptop & Rust (keep) \\
Inference engine & llama.cpp --- local mode, NOT RPC (keep, change usage) \\
Model & Meta Llama 3.2 1B/3B, GGUF, quantized (add) \\
Networking & TCP/JSON then gRPC over Wi-Fi (add) \\
Model split & by layers (pipeline) (add) \\
Recovery method & request-level re-route, controller-driven (core) \\
Code written with & Kiro \\
\bottomrule
\end{tabularx}
\end{table}

% ---------------------------------------------------------------------
\section{The Hardware}
\label{sec:hardware}
% ---------------------------------------------------------------------

\subsection{The five laptops (verified with the owner)}
\begin{table}[h]
\centering\footnotesize
\caption{Available demo hardware}
\label{tab:hardware}
\begin{tabularx}{\linewidth}{L{0.06\linewidth}L{0.16\linewidth}L{0.18\linewidth}L{0.14\linewidth}Y}
\toprule
\# & Laptop & GPU / chip & OS & llama.cpp backend \\
\midrule
1 & ASUS TUF & RTX 4050 & Windows & CUDA \\
2 & ASUS TUF & RTX 4050 & Windows & CUDA \\
3 & ASUS ROG & RTX 3050 & Linux & CUDA \\
4 & MacBook & Apple M4 & macOS & Metal \\
5 & MacBook & Apple M5 & macOS & Metal \\
\bottomrule
\end{tabularx}
\end{table}

Five laptops in total. Only three are needed for a convincing ``kill one, others
survive'' demo, so there is a comfortable margin.

\subsection{Why this mix helps}
\begin{itemize}[leftmargin=1.4em,nosep]
  \item \textbf{Two identical TUF 4050s} are the key asset: one can be a
  \textbf{perfect hot standby} for the other. Killing one 4050 and having its
  identical twin take the exact slice is the cleanest, most reliable failover
  demo possible.
  \item \textbf{Heterogeneous + cross-OS} (NVIDIA/Windows+Linux vs Apple/macOS)
  strengthens the story: ``different cheap laptops, different chips, different
  OSes, on Wi-Fi.'' Note: heterogeneity itself is \emph{not} the novelty
  (prima.cpp owns that, \VER); it is a bonus that makes the demo look real.
\end{itemize}

\subsection{The honest warning}
Even the ``core'' three NVIDIA machines are two OSes (2$\times$Windows + 1$\times$Linux),
and the Macs add a second backend (Metal). The Rust worker must build on
Windows, Linux, and macOS, and llama.cpp needs a CUDA build and a Metal build.
This is doable but is where early build time goes. The identical Windows 4050s
mean once the worker runs on one, the second is a copy. \HYP

\subsection{Recommended configuration}
\begin{itemize}[leftmargin=1.4em,nosep]
  \item \textbf{Core demo (must work), all-CUDA:} TUF 4050 \#1 (worker) + TUF
  4050 \#2 (worker + hot standby) + ROG 3050 (weaker worker, fewer layers,
  shows capability-aware assignment). One backend, one build path, lowest risk.
  \item \textbf{Stretch:} add M4 / M5 (Metal) as workers for a cross-OS ``pool''
  finale, and/or run the Go controller on a Mac.
  \item \textbf{Build order rule:} get the all-NVIDIA 3-node failover rock-solid
  \emph{before} touching CUDA$+$Metal mixing.
\end{itemize}

% ---------------------------------------------------------------------
\section{How It Works (Architecture)}
\label{sec:arch}
% ---------------------------------------------------------------------

\subsection{The normal path}
\begin{enumerate}[leftmargin=1.6em,nosep]
  \item The user sends a prompt to the \textbf{Go controller} over HTTP.
  \item The controller knows which worker holds which layers and sends the
  prompt down the pipeline.
  \item Worker~1 runs its layers, passes activations to worker~2, and so on.
  \item The last worker produces tokens, which stream back to the user.
  \item Throughout, every worker sends a \textbf{heartbeat} to the controller.
\end{enumerate}

\subsection{The failover path (the novel core)}
\begin{enumerate}[leftmargin=1.6em,nosep]
  \item A worker stops sending heartbeats (lid closed / Wi-Fi dropped / killed).
  \item The controller detects the gap within a short timeout.
  \item It picks a healthy machine that holds (or can load) the dead worker's
  layer slice --- preferring the identical hot standby.
  \item It reassigns the slice and \textbf{re-runs the current request}.
  \item Token streaming resumes; the user sees a short pause, not a crash.
\end{enumerate}

\subsection{Text diagram}
\begin{lstlisting}[style=plain]
                 +-------------------+
   user  ---->   |  Go CONTROLLER    |   heartbeats <---- all workers
   prompt        |  - task queue     |   failover logic lives here
                 |  - who-has-what   |
                 +---------+---------+
                           | assign layer slices
         +-----------------+-----------------+
         v                 v                 v
   +-----------+     +-----------+     +-----------+
   | WORKER 1  | --> | WORKER 2  | --> | WORKER 3  | --> tokens
   | 4050 (win)|     | 4050 (win)|     | 3050 (lin)|
   | layers0-8 |     | layers9-16|     |layers17-24|
   +-----------+     +-----------+     +-----------+
                      (hot standby
                       for worker 1)

   If WORKER 1 dies: controller detects silence ->
   reassigns layers 0-8 to the standby -> re-runs request ->
   tokens keep streaming.
\end{lstlisting}

\subsection{Why the controller, not llama.cpp, owns recovery}
llama.cpp's RPC mode aborts the process on any peer failure (\VER,
\S\ref{sec:landscape}). By keeping llama.cpp \emph{local} on each worker and
putting all coordination in the Go controller, a dead worker becomes a
controller event to handle --- not a fatal crash of the whole system.

% ---------------------------------------------------------------------
\section{The MVP Plan (Phases)}
\label{sec:mvp}
% ---------------------------------------------------------------------

The goal is one thing: \textbf{a worker can die mid-generation and the token
stream recovers}. Everything is sequenced to de-risk that, lowest-risk first.

\begin{table}[h]
\centering\footnotesize
\caption{Build phases}
\label{tab:phases}
\begin{tabularx}{\linewidth}{L{0.05\linewidth}L{0.26\linewidth}Y}
\toprule
\# & Phase & Goal / exit condition \\
\midrule
1 & One worker works & Rust worker loads a Llama GGUF via llama.cpp, runs its
layers, talks to the Go controller. Single-node path correct end to end. \\
2 & Two 4050s, static pipeline & Split a small Llama across the two identical
Windows 4050s over Wi-Fi. Correct tokens out. No failover yet. \\
3 & Failover (the core) & Heartbeat detects a dead 4050 → reassign its layers to
the twin → re-run request → tokens continue. Rock-solid on the identical pair. \\
4 & Add the Linux ROG 3050 & Third worker, weaker, gets fewer layers → shows
capability-aware assignment. Still CUDA, low extra risk. \\
5 & (Stretch) add Macs M4/M5 & Metal backend; cross-OS, cross-backend pool.
Only if phases 1--4 are done and stable. \\
6 & Demo + submission & Script the ``kill a laptop'' moment; \textbf{record a
backup video}; write the submission. \\
\bottomrule
\end{tabularx}
\end{table}

\subsection{Scope discipline}
\begin{itemize}[leftmargin=1.4em,nosep]
  \item \textbf{Validate at Phase~1} before investing in failover. If the
  single-node path is wrong, nothing downstream matters.
  \item \textbf{Macs are optional.} The submission must not depend on the
  trickiest (CUDA$+$Metal) integration.
  \item \textbf{Out of scope for MVP:} fine-tuning, security/auth, 70B models,
  phones, cross-Internet, mid-token handoff. List them as ``future work''.
\end{itemize}

\subsection{The must-do verification before committing (de-risking)}
\begin{enumerate}[leftmargin=1.6em,nosep]
  \item \textbf{Kill-test llama.cpp RPC} yourself: confirm it really dies on
  node loss in your setup (it should, per \VER). A $\sim$1-hour test that
  de-risks the whole premise.
  \item \textbf{Lock the exact Meta Llama GGUF} so the submission unambiguously
  ``uses a Meta model''.
  \item \textbf{Read the Official Rules} when Meta publishes them (promised
  before Oct~26) --- model-usage requirements can be stricter than the overview.
\end{enumerate}

% ---------------------------------------------------------------------
\section{Benchmarks (Measurable \& Honest)}
\label{sec:benchmarks}
% ---------------------------------------------------------------------

Measure the thing that is the differentiator --- recovery --- and be honest
about the thing that is not (raw speed).

\subsection{Primary metrics (the differentiator)}
\begin{table}[h]
\centering\footnotesize
\caption{Primary metrics}
\label{tab:primary}
\begin{tabularx}{\linewidth}{L{0.30\linewidth}Y L{0.18\linewidth}}
\toprule
Metric & Definition & Target \\
\midrule
Recovery time & Wall-clock from node-death to resumed token emission &
$<$ 2 seconds \\
Stream survival & \% of generations that complete correctly when 1 of N workers
is killed at a random token & $\sim$100\% for 1 failure \\
Output correctness & Tokens identical to a single-node greedy baseline (proves
failover does not corrupt output) & exact match \\
\bottomrule
\end{tabularx}
\end{table}

\subsection{Secondary metrics (context only --- where we do NOT win)}
\begin{itemize}[leftmargin=1.4em,nosep]
  \item Tokens/sec and time-per-output-token (TPOT) vs a single-laptop llama.cpp
  baseline and vs llama.cpp RPC multi-node (the closest apples-to-apples for LAN).
  \item \textbf{Do not stage a speed contest against prima.cpp} --- you will lose
  (\VER: prima.cpp reports 5--17$\times$ lower TPOT than llama.cpp/exo/dllama).
  Frame speed as ``acceptable overhead for resilience''.
\end{itemize}

\subsection{Baselines to run yourself}
\begin{enumerate}[leftmargin=1.6em,nosep]
  \item Single-node llama.cpp (the floor).
  \item llama.cpp RPC multi-node (to show it crashes on node kill --- the
  before/after contrast that sells the project).
  \item HyperPot (same kill, but survives).
\end{enumerate}
Context figures from prima.cpp (674~ms/token for 70B; 26~tok/s for 32B with
speculative decoding) are reference only. \VER

% ---------------------------------------------------------------------
\section{Hackathon Strategy (Meta Global AI Developer Hackathon)}
\label{sec:hackathon}
% ---------------------------------------------------------------------

\subsection{The verified facts about the event}
\begin{itemize}[leftmargin=1.4em,nosep]
  \item \VER\ \textbf{Dates:} build Oct~26 -- Nov~4 2026; submissions due Nov~4,
  9:00~PM ET; judging Nov~4--8; finals Nov~9.
  \item \VER\ \textbf{Track:} Indie developer track --- judged on
  \emph{creativity, technical ambition, and whether it works}.
  \item \VER\ \textbf{Models:} built around Meta Model API (Muse Spark, Muse
  Image, SAM 3.1), Meta open-weights models, and Muse Code; \$150 Meta Model API
  credits per participant.
  \item \VER\ \textbf{Themes:} ``Build for Your World'' or ``Push the Frontier''.
\end{itemize}

\subsection{Why Indie fits HyperPot}
The Indie criterion ``whether it works'' is the purest possible match for a live
demo where a laptop is physically killed and the AI keeps talking. Technical
ambition is covered by bringing datacenter-grade recovery to cheap laptops.

\subsection{The pitch (safe wording)}
\begin{quote}
``Fault tolerance like this exists in datacenters. It has never run on a pile of
cheap student laptops over Wi-Fi with free tools. We made that work --- watch,
I'll kill one.''
\end{quote}
This is defensible because every clause is backed by a source
(\S\ref{sec:landscape}). \textbf{Never claim} to have invented the algorithm ---
LUMEN and others got there first (\VER); claim that you \emph{brought it down to
laptops}, which is true and humble.

\subsection{Meta-model requirement}
Run a Meta open-weights Llama (small GGUF) through the pipeline. This satisfies
``use a Meta model''. ``Push the Frontier'' is the natural theme (a new
capability on commodity hardware), though ``Build for Your World'' works if
wrapped around a visible app.

\subsection{Demo safety}
Always have a \textbf{recorded backup video} of the kill-and-recover moment.
Never rely solely on live Wi-Fi in front of judges.

% ---------------------------------------------------------------------
\section{Risks \& Immediate Next Steps}
\label{sec:risks}
% ---------------------------------------------------------------------

\subsection{The honest risks}
\begin{itemize}[leftmargin=1.4em,nosep]
  \item \textbf{llama.cpp fights failover.} The abort on peer failure is baked
  into its RPC layer (\VER). Mitigation: use llama.cpp \emph{locally} and do
  recovery in the controller (\S\ref{sec:arch}); use request-level re-route, not
  mid-token handoff.
  \item \textbf{Cross-platform build cost.} Windows + Linux + macOS, CUDA +
  Metal. Mitigation: all-CUDA core first; Macs are stretch only. \HYP
  \item \textbf{Novelty challenge from an informed judge.} Mitigation: the
  ``brought it to laptops'' framing (\S\ref{sec:hackathon}), never ``invented''.
  \item \textbf{Live-demo failure.} Mitigation: recorded backup video.
\end{itemize}

\subsection{Guardrails being held}
\begin{itemize}[leftmargin=1.4em,nosep]
  \item \Keep: Go controller, Rust worker, llama.cpp (local), Meta Llama (small
  GGUF), Wi-Fi.
  \item \Drop: fine-tuning, llama.cpp RPC mode, 70B models, mid-token handoff.
  \item Win condition: a working live failover demo --- not a novelty claim.
\end{itemize}

\subsection{Immediate next steps}
\begin{enumerate}[leftmargin=1.6em,nosep]
  \item Kill-test llama.cpp RPC to confirm the premise (\S\ref{sec:mvp}).
  \item Lock the exact Meta Llama GGUF.
  \item Read the Official Rules once published.
  \item Build Phase~1 (one Rust worker + Go controller + llama.cpp, single node)
  --- only on an explicit ``build'' instruction.
\end{enumerate}

\subsection{Status of this document}
This is a planning and research artifact. No HyperPot code has been written yet.
The design extends the existing Go~+~Rust~+~llama.cpp architecture and is
sequenced so the riskiest piece (failover on identical hardware) is proven before
any cross-platform work begins.

% ---------------------------------------------------------------------
\section{Evidence Log: Kill-Test on Real Hardware}
\label{sec:evidence}
% ---------------------------------------------------------------------

This chapter records the de-risking experiment run \emph{before} any HyperPot
code was written, to verify the core premise on the actual demo laptops rather
than relying only on documentation.

\subsection{Setup}
\begin{table}[h]
\centering\footnotesize
\caption{Kill-test configuration (26 Oct 2026)}
\label{tab:killtest}
\begin{tabularx}{\linewidth}{L{0.26\linewidth}Y}
\toprule
Item & Value \\
\midrule
Laptop A (main process) & ROG, RTX 3050 Laptop GPU (3767 MiB), Manjaro Linux, CUDA 13.3 \\
Laptop B (remote node) & TUF, RTX 4050, Windows, llama.cpp b11900 CPU build \\
Engine & llama.cpp v0.7.0-dev (commit caa8191), built with \code{-DGGML\_CUDA=ON -DGGML\_RPC=ON} on the ROG \\
Model & Llama-3.2-1B-Instruct, Q4\_K\_M GGUF (771 MB) \\
Link & Wi-Fi; B exposed \code{ggml-rpc-server -H 0.0.0.0 -p 50052} \\
Connectivity & TCP port 50052 reachable from A (verified) \\
\bottomrule
\end{tabularx}
\end{table}

\subsection{Run 1 --- normal distributed inference (baseline)}
A full generation ran across A~+~B over RPC and \textbf{completed successfully},
producing a coherent multi-paragraph story. Final stats:
\code{[ Prompt: 52.7 t/s | Generation: 7.7 t/s ]}, then \code{Exiting...}.
\textbf{Conclusion:} distributed inference across two heterogeneous laptops over
Wi-Fi \emph{works}. \VER

\subsection{Run 2 --- remote node lost (the premise)}
A longer run was started. During model data transfer to the remote node, the
Wi-Fi link to B dropped packets and the remote node became unreachable. The main
process on A did \textbf{not} recover --- it aborted the entire process:
\begin{lstlisting}[style=plain]
ggml/src/ggml-rpc/ggml-rpc.cpp:685: Remote RPC server crashed
    or returned malformed response
E send failed (bytes_sent=6877944, size_to_send=2559545)
Aborted (core dumped)
\end{lstlisting}

\subsection{What this proves}
\begin{itemize}[leftmargin=1.4em,nosep]
  \item \VER\ \textbf{llama.cpp RPC has no fault tolerance on this exact
  hardware.} Loss of the remote node terminates the whole process
  (\code{ggml-rpc.cpp:685} $\rightarrow$ \code{Aborted}), matching the
  documentation (\S\ref{sec:landscape}).
  \item \VER\ \textbf{The failure is realistic, not contrived.} It was triggered
  by ordinary Wi-Fi instability --- no deliberate kill was even needed. This is
  precisely the scenario HyperPot targets.
  \item \textbf{Answer to ``does it restart from the beginning?''} No --- plain
  llama.cpp does not restart at all; it crashes and must be re-launched by hand.
  HyperPot's contribution is to make the controller detect the loss, re-route the
  work, and \textbf{automatically re-run the request} so generation continues.
\end{itemize}

\subsection{Note for the live demo}
The crash here came partly from random Wi-Fi drops during a large transfer to a
CPU-only node. For a reproducible stage demo, the ``kill'' should be a
\textbf{deliberate, clean action} (Ctrl+C on the remote server, or lid-close)
rather than relying on chance packet loss. The before/after contrast ---
llama.cpp aborts vs HyperPot recovers --- is the core of the demonstration. Both
log files (\code{killtest.log}, \code{killtest2.log}) are retained as evidence.

% ---------------------------------------------------------------------
\section{References \& Sources}
\label{sec:references}
% ---------------------------------------------------------------------

Every source below was consulted directly during the research for this document.
Each is tagged by how it was used. Access dates: October 2026.

\sloppy

\subsection{Research papers (competitive landscape)}
\begin{enumerate}[leftmargin=1.8em,label={[P\arabic*]}]
  \item \textbf{Prima.cpp: Fast 30--70B LLM Inference on Heterogeneous and
  Low-Resource Home Clusters.} Z.~Li et al., arXiv:2504.08791 (v1 Apr 2025, v3
  Jul 2026). \url{https://arxiv.org/abs/2504.08791}. \\
  \emph{Used for:} heterogeneous/Wi-Fi home-cluster inference; Halda scheduler;
  the 5--17$\times$ lower TPOT claim vs llama.cpp/exo/dllama; static scheduling
  (no mid-inference recovery). \VER

  \item \textbf{Petals: Collaborative Inference and Fine-tuning of Large Models.}
  A.~Borzunov et al., arXiv:2209.01188. \url{https://arxiv.org/abs/2209.01188}. \\
  \emph{Used for:} BitTorrent-style swarm; exposing hidden states for fine-tuning. \VER

  \item \textbf{Distributed Inference and Fine-tuning of Large Language Models
  Over The Internet.} A.~Borzunov et al., arXiv:2312.08361. \\
  \url{https://arxiv.org/abs/2312.08361}. \\
  \emph{Used for:} churn-tolerant inference over the Internet for 50B+ models
  (``any device can disconnect abruptly''). \VER

  \item \textbf{Distributed Llama --- Distributed Inference of Large Language
  Models with Slow Synchronization over Ethernet.} B.~Tadych, 2024. \\
  \url{https://github.com/b4rtaz/distributed-llama} (report.pdf). \\
  \emph{Used for:} tensor-parallel home-device inference; synchronization cost. \VER

  \item \textbf{LUMEN: Coordinated Failure Recovery for Distributed LLM Serving.}
  Z.~Cao et al., arXiv:2606.17787. \url{https://arxiv.org/abs/2606.17787}. \\
  \emph{Used for:} proof that load-aware, interrupted-request recovery
  \emph{already exists} in research (datacenter/GPU). Key to the honest-novelty
  statement. \VER

  \item \textbf{Intelligent Orchestration of Distributed Large Foundation Model
  Inference at the Edge.} arXiv:2504.03668. \\
  \url{https://arxiv.org/abs/2504.03668}. \\
  \emph{Used for:} capacity-aware workload distribution + dynamic partition
  migration + real-time re-split (the capability-aware recovery idea, already
  published). \VER

  \item \textbf{FailSafe: High-performance Resilient Serving.} arXiv:2511.14116.
  \url{https://arxiv.org/abs/2511.14116}. \\
  \emph{Used for:} serving that survives multiple GPU failures. \VER

  \item \textbf{PipeLive: Efficient Live In-place Pipeline Parallelism
  Reconfiguration for Dynamic LLM Serving.} arXiv:2604.12171. \\
  \url{https://arxiv.org/abs/2604.12171}. \\
  \emph{Used for:} live pipeline reconfiguration during serving. \VER

  \item \textbf{DynaPipe: Dynamic Layer Redistribution for Pipeline-Parallel LLM
  Inference.} NeurIPS 2025. \\
  \emph{Used for:} dynamic layer redistribution + async KV-cache migration during
  inference. \VER

  \item \textbf{EnerInfer: Energy-Aware On-Device LLM Inference.}
  arXiv:2606.23001. \url{https://arxiv.org/abs/2606.23001}. \\
  \emph{Used for:} evidence that battery/thermal-aware on-device LLM scheduling
  \emph{already exists} --- which ruled out that angle as a novelty. \VER

  \item \textbf{LLM Inference at the Edge: Mobile, NPU, and GPU Performance
  Efficiency Trade-offs Under Sustained Load.} arXiv:2603.23640. \\
  \emph{Used for:} thermal throttling as the dominant constraint on-device. \VER

  \item \textbf{MNN-AECS: Energy Optimization for LLM Decoding on Mobile Devices
  via Adaptive Core Selection.} arXiv:2506.19884. \\
  \emph{Used for:} further prior art on energy-adaptive on-device decoding. \VER
\end{enumerate}

\subsection{Official repositories \& documentation}
\begin{enumerate}[leftmargin=1.8em,label={[R\arabic*]}]
  \item \textbf{llama.cpp --- RPC tool README.} ggml-org/llama.cpp. \\
  \url{https://github.com/ggml-org/llama.cpp/blob/master/tools/rpc/README.md}. \\
  \emph{Used for:} the RPC backend design and its insecurity warning. \VER

  \item \textbf{llama.cpp RPC --- peer-failure behaviour.} LlamaCppEx binding
  docs. \url{https://llama-cpp-ex.hexdocs.pm/LlamaCppEx.RPC.html}. \\
  \emph{Used for:} the verified statement that a peer crash/drop terminates the
  process (\code{RPC\_STATUS\_ASSERT}/\code{GGML\_ABORT}). \VER

  \item \textbf{llama.cpp issue \#11508} --- crash while loading in multi-node
  RPC after a node crash. \url{https://github.com/ggml-org/llama.cpp/issues/11508}. \VER

  \item \textbf{llama.cpp issue \#9214} --- RPC backend crash. \\
  \url{https://github.com/ggml-org/llama.cpp/issues/9214}. \VER

  \item \textbf{Exo.} exo-explore/exo. \url{https://github.com/exo-explore/exo}. \\
  \emph{Used for:} features (auto-discovery, MLX, Thunderbolt RDMA, tensor
  parallel); open issues on instability (\#1726, \#1847, \#2007, \#775). \VER

  \item \textbf{Deep Dive: Exo --- Distributed AI Inference on Consumer Hardware}
  (review, Sept 2025). \emph{Used for:} ``critical gaps in security, fault
  tolerance, and operational tooling''. \VER

  \item \textbf{prima.cpp repository.} \url{https://github.com/OpenCPIL/prima.cpp}
  (also gitee.com/zonghang-li/prima.cpp). \VER

  \item \textbf{QVAC Fabric LLM} (Tether Data). \\
  \url{https://qvac.tether.io/products/fabric} and the Dec~2025 launch
  announcement on \code{tether.io}. \\
  \emph{Used for:} existence of an edge-first inference + LoRA fine-tuning
  framework. Existence \VER; technical internals \UNK.
\end{enumerate}

\subsection{Hackathon (rules \& requirements)}
\begin{enumerate}[leftmargin=1.8em,label={[H\arabic*]}]
  \item \textbf{Meta Global AI Developer Hackathon --- official page.} \\
  \url{https://dev.meta.ai/events/global-hackathon}. \\
  \emph{Used for:} dates, Indie vs Startup tracks and their judging criteria,
  Meta Model API / open-weights model usage, \$150 credits, themes, eligibility. \VER

  \item \textbf{Cerebral Valley event listing.} \\
  \url{https://cerebralvalley.ai/e/meta-ai-developer-hackathon}. \\
  \emph{Used for:} the 20,000-applicant cap and rolling evaluation. \VER
\end{enumerate}

\subsection{Evidence-strength note}
Items marked \VER\ are quoted or closely paraphrased from the named source.
Where a source is marketing material (e.g. QVAC Fabric), only the
\emph{existence} of the capability is treated as verified; its internal
mechanisms remain \UNK\ until primary technical documentation is available. No
claim of HyperPot's novelty rests on an \UNK\ item.


\end{document}
